%% Document 168 – Chapter (English version)
%% The Periodic Table as xi-Geometry:
%% Radial Mass Ladder and Angular Shell Structure from a Common Origin
%% Johann Pascher, March 2026

\chapter{The Periodic Table as \texorpdfstring{$\xi$}{xi}-Geometry}
\label{chap:periodic-xi}

\begin{abstract}
The periodic table of the elements and the particle-mass ladder of T0 theory
share the same geometric quantum number~$n$, but in different roles.
Particle masses follow a \emph{radial} exponential scaling
$m_n \sim \xi^{-n} \cdot m_P$, while the shell structure of the periodic table
obeys an \emph{angular} quadratic degeneracy $g_n = 2n^2$.
Both arise from the three-dimensional lattice geometry of T0 theory:
$\xi$ is the packing deficit of the 3D lattice (radial direction),
$2n^2$ is the degeneracy of the $n$-th spherical shell in the same lattice
(angular direction).
The Rydberg energy lies exactly on rung~7 of the $\xi$-power ladder:
$R_\infty = \frac{m_e \alpha^2}{2} \approx \xi^{6.956} \cdot E_P$.
The duality principle $E_n \cdot g_n = 2R_\infty = \text{const}$
connects the $1/n^2$ energy ladder with the $n^2$ filling ladder
and is a geometric equipartition.
The periodic table is the angular cross-section of the $\xi$-geometry at
rung $\xi^{\sim 7}$ — analogously, the lepton masses are the radial cross-section
at rungs $\xi^{4.9}$–$\xi^{5.8}$.
\end{abstract}

%──────────────────────────────────────────────────────────────────────────────
\section{Introduction and Problem Statement}
\label{sec:intro}

T0 theory (Time-Mass Duality / FFGFT) derives all physical constants from the
geometric parameter
\begin{equation}
  \xi = \frac{4}{30\,000} = 1.\overline{3} \times 10^{-4}
\end{equation}
\cite{dok009,dok041,dok056}.
The particle masses of the Standard Model organise themselves on an exponential
ladder
\begin{equation}
  m_n \;\sim\; \xi^{-n} \cdot m_P, \qquad n = 1, 2, 3, \ldots
  \label{eq:massldr}
\end{equation}
where $n$ is the T0 quantum number and $m_P$ is the Planck mass.

The periodic table of the elements, on the other hand, shows a very different
pattern: shell capacities follow
\begin{equation}
  g_n \;=\; 2n^2 \;=\; 2,\; 8,\; 18,\; 32,\; 50,\;\ldots
  \label{eq:shells}
\end{equation}
and binding energies scale as $1/n^2$.
At first glance, particle masses and the periodic table appear unrelated.

\textbf{Gemini's observation} (2026): the periodic table lies ``one level
higher'' than the particle masses — both have quantum numbers, but what
common geometric origin connects their patterns?

The answer of this document: \emph{T0 geometry has two axes —
a radial one ($\xi^n$ scaling) and an angular one ($2n^2$ degeneracy),
both arising from the three-dimensional lattice structure with packing deficit~$\xi$.}

%──────────────────────────────────────────────────────────────────────────────
\section{The \texorpdfstring{$\xi$}{xi}-Power Ladder and Its Rungs}
\label{sec:xi-ladder}

Previous documents in the T0 series have built up the $\xi$-power ladder
step by step. Table~\ref{tab:xi-ladder} shows the current state including
the new atomic rung.

\begin{table}[htbp]
\centering
\caption{Complete $\xi$-power ladder (as of March 2026)}
\label{tab:xi-ladder}
\begin{tabular}{clll}
\hline
Rung & $\xi$ power & Physical content & Document \\
\hline
$0$      & $\xi^0 = 1$                 & Planck scale, $L_0 = \xi \cdot \ell_P$ & Doc.\,009 \\
$1/2$    & $\xi^{1/2} \approx 10^{-2}$ & $\alpha_W$, $m_e/m_\mu$, PPLN $\Delta n$ & Doc.\,041, 167 \\
$1$      & $\xi^1 \approx 10^{-4}$     & $\alpha_\text{EM} \sim \xi \cdot E_0^2$ & Doc.\,011 \\
$4$      & $\xi^4 \approx 10^{-16}$    & CMB temperature $T_\text{CMB}$           & Doc.\,166 \\
$23/4$   & $\xi^{5.775}$               & Electron mass $m_e/m_P$                  & Doc.\,056 \\
$\mathbf{\sim 7}$ & $\boldsymbol{\xi^{6.956}}$ & \textbf{Rydberg energy $R_\infty$, periodic table} & \textbf{this doc.} \\
$10$     & $\xi^{10} \approx 10^{-43}$ & Hubble constant $H_0$                    & Doc.\,165 \\
\hline
\end{tabular}
\end{table}

The Rydberg energy lies between rungs~6 and~7:
\begin{equation}
  \frac{R_\infty}{E_P}
  = \frac{m_e}{m_P} \cdot \frac{\alpha^2}{2}
  = \xi^{5.775} \cdot \xi^{1.181}
  = \xi^{6.956}
  \approx \xi^7
\end{equation}
The exponent is the sum of two independent T0 quantities: the electron mass
(rung $5.775$, Doc.\,056) and the fine-structure constant (rung $1.181$,
Doc.\,011).

%──────────────────────────────────────────────────────────────────────────────
\section{Two Axes of \texorpdfstring{$\xi$}{xi}-Geometry}
\label{sec:two-axes}

\subsection{Radial Axis: Exponential Mass Ladder}
\label{subsec:radial}

Particle masses organise themselves along the \emph{radial} axis of the T0
lattice. The fundamental relation is \cite{dok009}:
\begin{equation}
  m_n \;\sim\; \xi^{-n} \cdot m_P, \qquad n = 1, 2, 3, \ldots
\end{equation}
Each step on this ladder multiplies the mass by $\xi^{-1} = 7500$.
The leptons occupy the rungs (Doc.\,056, 122):
\begin{align}
  m_e:    &\quad \xi^{5.775} \cdot m_P \\
  m_\mu:  &\quad \xi^{5.177} \cdot m_P \\
  m_\tau: &\quad \xi^{4.861} \cdot m_P
\end{align}
The mass ratio $m_e/m_\mu \propto \xi^{1/2}$ (Doc.\,041) lies exactly on the
electroweak exponent — the half-rung of the $\xi$-ladder.

\subsection{Angular Axis: Quadratic Degeneracy Ladder}
\label{subsec:angular}

The \emph{angular} axis of the same lattice determines the degeneracy of
spherical shells. In 3D space, the $n$-th principal shell has:
\begin{equation}
  g_n \;=\; 2 \sum_{l=0}^{n-1} (2l+1) \;=\; 2n^2
  \label{eq:gn}
\end{equation}
The factor 2 comes from spin $m_s = \pm\tfrac{1}{2}$; the sum is pure
angular-momentum algebra in 3D space. Table~\ref{tab:shells} shows the
first four shells.

\begin{table}[htbp]
\centering
\caption{Shell structure and the periodic table}
\label{tab:shells}
\begin{tabular}{ccccc}
\hline
$n$ & $g_n = 2n^2$ & Period length & Noble gas $Z$ & Period number \\
\hline
1 &  2 &  2 &  He (2)  & 1 \\
2 &  8 &  8 &  Ne (10) & 2, 3 \\
3 & 18 & 18 &  Ar (18) & 4, 5 \\
4 & 32 & 32 &  Kr (36) & 6, 7 \\
\hline
\end{tabular}
\end{table}

\subsection{The Connection through Quantum Number \texorpdfstring{$n$}{n}}
\label{subsec:connection-n}

The same quantum number~$n$ appears in both systems, but with different
functional roles:
\begin{align}
  \text{Radial (masses):}  &\quad m_n \;\sim\; \xi^{-n} \cdot m_P
    \quad\text{(exponential)} \\
  \text{Angular (shells):} &\quad g_n \;=\; 2n^2
    \quad\text{(quadratic)}
\end{align}

This is not coincidental. In T0 lattice geometry, $n$ corresponds to the
\emph{discrete radius in the lattice}:
\begin{itemize}
  \item Radially: $n$ steps outward from the lattice $\Rightarrow$
    mass scales as $\xi^{-n}$ (geometric sequence)
  \item Angularly: the lattice points on the sphere of radius~$n$ carry
    degeneracy $2n^2$ (combinatorial sequence)
\end{itemize}

%──────────────────────────────────────────────────────────────────────────────
\section{The Duality Principle}
\label{sec:duality}

\subsection{Energy Ladder and Filling Ladder Are Dual}
\label{subsec:dual}

On the atomic scale (hydrogen-like atoms):
\begin{align}
  \text{Energy ladder:}  &\quad E_n \;=\; -\frac{R_\infty}{n^2}
    \quad\text{(falls as } 1/n^2\text{)} \\
  \text{Filling ladder:} &\quad g_n \;=\; 2n^2
    \quad\text{(grows as } n^2\text{)}
\end{align}

Their product is exactly constant:
\begin{equation}
  \boxed{E_n \cdot g_n \;=\; \frac{R_\infty}{n^2} \cdot 2n^2 \;=\; 2R_\infty
  \;=\; \text{const}}
  \label{eq:duality}
\end{equation}

This is geometric \emph{equipartition}: every shell contributes the same total
energy $2R_\infty$ to the Rydberg scale, regardless of~$n$.
The $1/n^2$ energy ladder and the $n^2$ filling ladder are \emph{inversely dual}.

\subsection{Comparison with the Mass Ladder}

Table~\ref{tab:compare} summarises all three ladders.

\begin{table}[htbp]
\centering
\caption{Three ladders — one geometric origin}
\label{tab:compare}
\begin{tabular}{lccc}
\hline
System & Ladder & Function of $n$ & Type \\
\hline
Particle masses & $m_n \sim \xi^{-n} \cdot m_P$  & exponential & radial \\
Shell energy    & $E_n = R_\infty / n^2$          & harmonic    & angular \\
Shell filling   & $g_n = 2n^2$                    & quadratic   & angular \\
\hline
\multicolumn{4}{l}{Duality: $E_n \cdot g_n = 2R_\infty = \text{const}$} \\
\hline
\end{tabular}
\end{table}

%──────────────────────────────────────────────────────────────────────────────
\section{The Role of Fractal Dimension}
\label{sec:fractal}

In T0 (Doc.\,009, 133):
\begin{equation}
  D_f \;=\; 3 - \xi \;=\; 2.999867
\end{equation}
Space dimension is almost, but not exactly, 3.

In the limit $D_f \to 3$ ($\xi \to 0$):
\begin{itemize}
  \item $g_n = 2n^2$ is \emph{exact} (pure spherical symmetry in the $\mathbb{Z}^3$ lattice)
  \item the shell structure of the periodic table is exact
\end{itemize}

For $D_f = 3 - \xi < 3$ there is a minimal correction:
\begin{equation}
  \delta g_n \;\sim\; \xi \cdot n^2
\end{equation}
For $n = 2$: $\delta g_2 \approx \xi \cdot 4 \approx 5 \times 10^{-4}$
(entirely negligible).

The periodic table is therefore exact because $D_f \approx 3$ to four decimal
places. The tiny deviation $\xi \approx 1.3 \times 10^{-4}$ makes space
``almost-integer-dimensional'' — and precisely this deviation generates all
the physics of particle masses and fundamental constants.

%──────────────────────────────────────────────────────────────────────────────
\section{Numerical Verification}
\label{sec:numerics}

\subsection{Rydberg Energy on the \texorpdfstring{$\xi$}{xi}-Ladder}

\begin{align}
  R_\infty &= \frac{m_e \alpha^2}{2} = 13.5984\;\text{eV} \\
  \frac{R_\infty}{E_P} &= \frac{m_e}{m_P} \cdot \frac{\alpha^2}{2}
    = \xi^{5.775} \cdot \xi^{1.181} = \xi^{6.956}
\end{align}

\begin{table}[htbp]
\centering
\caption{Rydberg energy: T0 derivation vs. measurement}
\label{tab:rydberg}
\begin{tabular}{lcc}
\hline
Quantity & Value & $\xi$-exponent \\
\hline
$m_e / m_P$        & $4.185 \times 10^{-23}$ & $5.775$ \\
$\alpha^2 / 2$     & $2.663 \times 10^{-5}$  & $1.181$ \\
$R_\infty / E_P$   & $1.114 \times 10^{-27}$ & $6.956$ \\
$R_\infty$ (T0, approx.) & $\approx 9.1\;\text{eV}$ ($\xi^7$) & $7.000$ \\
$R_\infty$ (measured)    & $13.598\;\text{eV}$ & $6.956$ \\
\hline
\end{tabular}
\end{table}

\subsection{Duality Relation: Numerical Verification}

\begin{align}
  E_1 \cdot g_1 &= 13.60\;\text{eV} \cdot 2  = 27.20\;\text{eV} = 2R_\infty \\
  E_2 \cdot g_2 &= \phantom{0}3.40\;\text{eV} \cdot 8  = 27.20\;\text{eV} = 2R_\infty \\
  E_3 \cdot g_3 &= \phantom{0}1.51\;\text{eV} \cdot 18 = 27.18\;\text{eV} \approx 2R_\infty \\
  E_4 \cdot g_4 &= \phantom{0}0.85\;\text{eV} \cdot 32 = 27.20\;\text{eV} = 2R_\infty
\end{align}

The duality relation~\eqref{eq:duality} holds exactly for all~$n$.

%──────────────────────────────────────────────────────────────────────────────
\section{The Periodic Table in the T0 Hierarchy}
\label{sec:hierarchy}

\begin{figure}[htbp]
\centering
\begin{tabular}{|l|l|l|}
\hline
\textbf{$\xi$-rung} & \textbf{Physics} & \textbf{Axis} \\
\hline
$\xi^0 = 1$   & Planck scale ($L_0 = \xi \cdot \ell_P$)    & — \\
$\xi^{1/2}$   & Electroweak: $\alpha_W$, $m_e/m_\mu$         & radial \\
$\xi^1$       & $\alpha_\text{EM}$, EM coupling               & radial \\
$\xi^4$       & CMB temperature $T_\text{CMB}$                & radial \\
$\xi^{5.8}$   & Particle masses (leptons, quarks)             & radial \\
$\xi^{6.956}$ & \textbf{Rydberg $R_\infty$, periodic table}  & \textbf{both} \\
$\xi^{10}$    & Hubble constant $H_0$                         & radial \\
\hline
\end{tabular}
\caption{Position of the periodic table in the complete $\xi$-hierarchy}
\label{fig:hierarchy}
\end{figure}

\vspace{0.5em}

The special feature of the atomic rung $\xi^{\sim 7}$: here \emph{both} axes
act. The radial axis determines the energy scale $R_\infty$; the angular axis
determines the pattern $2n^2$. The periodic table is the visible result of
both axes acting together.

%──────────────────────────────────────────────────────────────────────────────
\section{Predictions and Open Questions}
\label{sec:predictions}

\subsection{Angular Corrections from \texorpdfstring{$D_f < 3$}{Df<3}}

The fractal dimension $D_f = 3 - \xi$ leads to minimal corrections of the
exact $2n^2$ rule:
\begin{equation}
  g_n^\text{(T0)} \;=\; 2n^2 \cdot (1 - \xi \cdot f(n))
\end{equation}
where $f(n)$ is a dimensionless function of~$n$ (still to be derived).
For $n \leq 4$ the correction is $< 10^{-3}$ and not measurable.

\subsection{Molecular Binding Energies}

Chemical binding energies lie in the range $1$–$10\;\text{eV}$, hence also
on rung $\xi^7$. It remains to be checked whether specific binding energies
lie on rational multiples of $\xi^{7/2} \cdot R_\infty$.

\subsection{Heavy Elements and Relativistic Corrections}

For $Z \gg 1$ relativistic effects become important. In T0 these are connected
with the $\xi^{1/2}$ electroweak rung. A systematic analysis of ionisation
energies of heavy elements ($Z > 54$) might reveal $\xi$-matches at the
half-rung $\xi^{7-1/2} = \xi^{13/2}$.

%──────────────────────────────────────────────────────────────────────────────
\section{Summary}
\label{sec:summary}

\begin{enumerate}
  \item \textbf{Energy scale:} $R_\infty / E_P = \xi^{6.956} \approx \xi^7$.
    The Rydberg energy — and hence all of atomic physics — lies on rung~7
    of the $\xi$-power ladder.

  \item \textbf{Radial axis:} Particle masses follow $m_n \sim \xi^{-n} \cdot m_P$.
    The exponent~$n$ is the discrete lattice radius in the radial $\xi$-direction.

  \item \textbf{Angular axis:} Shell capacities follow $g_n = 2n^2$.
    The index~$n$ is the same discrete lattice radius, now acting as a shell
    index in the angular 3D-sphere direction.

  \item \textbf{Duality principle:} $E_n \cdot g_n = 2R_\infty = \text{const}$.
    The $1/n^2$ energy ladder and the $n^2$ filling ladder are inversely dual.
    Every shell contributes the same total energy $2R_\infty$ to the $\xi^7$ scale.

  \item \textbf{Fractal dimension:} $D_f = 3 - \xi \approx 3$ makes $2n^2$ exact
    (to $\sim 10^{-4}$). The periodic table is exact because space is
    almost-integer-dimensional.

  \item \textbf{Common origin:} Particle masses and the periodic table arise from
    the same 3D lattice geometry with packing deficit~$\xi$.
    The periodic table is the \emph{angular cross-section} of $\xi$-geometry
    at rung $\xi^{\sim 7}$; the lepton masses are the \emph{radial cross-section}
    at rungs $\xi^{4.9}$–$\xi^{5.8}$.
\end{enumerate}

\begin{thebibliography}{99}
\bibitem{dok009} Pascher, J.: \emph{Origin of the $\xi$ Parameter (Doc.\,009)},
  T0 Theory Series, 2025.
\bibitem{dok011} Pascher, J.: \emph{Fine-Structure Constant (Doc.\,011)},
  T0 Theory Series, 2025.
\bibitem{dok041} Pascher, J.: \emph{Parameter Derivation (Doc.\,041)},
  T0 Theory Series, 2025.
\bibitem{dok056} Pascher, J.: \emph{Universal Derivation of Lepton Masses (Doc.\,056)},
  T0 Theory Series, 2025.
\bibitem{dok122} Pascher, J.: \emph{$m_e/m_\mu$ Ratio (Doc.\,122)},
  T0 Theory Series, 2025.
\bibitem{dok133} Pascher, J.: \emph{Fractal Correction $K_\text{frak}$ (Doc.\,133)},
  T0 Theory Series, 2025.
\bibitem{dok165} Pascher, J.: \emph{Hubble Ratio on the $\xi$-Ladder (Doc.\,165)},
  T0 Theory Series, 2026.
\bibitem{dok166} Pascher, J.: \emph{Casimir-CMB-$H_0$ Bridge (Doc.\,166)},
  T0 Theory Series, 2026.
\bibitem{dok167} Pascher, J.: \emph{LiNbO$_3$ on the $\xi$-Ladder (Doc.\,167)},
  T0 Theory Series, 2026.
\end{thebibliography}
